\documentclass[11pt,reqno]{amsart}

\usepackage{amsmath,amssymb,amsthm}
\usepackage[T1]{fontenc}
\usepackage[margin=1in]{geometry}
\usepackage{hyperref}
\hypersetup{colorlinks=true,linkcolor=blue,citecolor=blue,urlcolor=blue}

\theoremstyle{plain}
\newtheorem{theorem}{Theorem}
\newtheorem{proposition}{Proposition}
\newtheorem{lemma}{Lemma}
\newtheorem{corollary}{Corollary}

\theoremstyle{definition}
\newtheorem{definition}{Definition}
\newtheorem{remark}{Remark}

\DeclareMathOperator{\rank}{rank}
\newcommand{\Q}{\mathbb{Q}}
\newcommand{\Z}{\mathbb{Z}}
\newcommand{\Pone}{\mathbb{P}^{1}}
\newcommand{\Epcp}{E_{\mathrm{PCP}}}

\begin{document}

\title[Degeneracy of torsion under the cuboid recovery map]{Torsion of the perfect-cuboid elliptic family is degenerate
for the cuboid recovery map}

\author{C$\Lambda$ / Lightman Chang}
\address{Independent Researcher}
\email{lightman.chang@gmail.com}

\date{\today}

\subjclass[2020]{Primary 11G05; Secondary 11D09, 11G07}

\keywords{Perfect cuboid, Euler brick, elliptic curve, rational torsion,
Mazur's theorem, recovery map, degeneracy}

\begin{abstract}
For a Pythagorean rational $q$ (one with $1+q^{2}$ a rational square) the
perfect-cuboid attack attaches to the corresponding fibre the elliptic curve
$\Epcp(q)\colon Y^{2}=X(X+1)(X+q^{2})$ together with a rational recovery map
$\varphi(X,Y)=2Yq/(q^{2}-X^{2})$ that returns the candidate cuboid edge of a rational
point. The rational torsion of this family is $\Z/4\Z\times\Z/2\Z$. This
classification is not new: in substance and method (Mazur's theorem followed by a Fermat
descent on $u^{2}=s^{4}+1$) it is that of Yoshida's Lemma~2.1 for his face-cuboid family
$E_{1,s}$, which we credit accordingly and re-derive through an explicit uniform
discriminant identity $\Delta(Z)=(Z-1)^{4}(Z+1)^{2}(Z^{2}-6Z+1)$; we also record that
$\Epcp(q)$ and $E_{1,s}$ are not $\Q$-isomorphic on the Pythagorean locus, only
twist-related. The contribution of this note is one statement, with no analogue in
Yoshida's work, on the interaction of the torsion with the recovery map: every
non-identity rational torsion point of $\Epcp(q)$ is sent by $\varphi$ into
$\{0,\infty\}$, the three two-torsion points to $0$ and the four order-four points to a
pole (the identity being the point at infinity). Hence a torsion
point never decodes to a finite nonzero edge, and any perfect-cuboid point on the fibre
is of infinite order. The degeneracy is proved as an identity in $\Q(q)$ and confirmed
numerically across $62$ Pythagorean values. We offer it as a companion structural remark
to Yoshida and as the lemma on which a Chabauty- or descent-type search of the
non-torsion locus would rest. No claim is made about the existence of perfect cuboids.
\end{abstract}

\maketitle

\section{Introduction}

A \emph{perfect cuboid} is a rectangular box whose three edges, three face diagonals,
and space diagonal are all rational; whether one exists is a long-standing open
question \cite{Guy}. A productive way to organise the search is to slice the
perfect-cuboid locus into one-parameter fibres indexed by a Pythagorean rational $q$
(a rational with $1+q^{2}$ a rational square) and to attach to each fibre an elliptic
curve whose rational points carry the arithmetic of that fibre. In the parametrisation
used here the relevant curve is
\begin{equation}\label{eq:Epcp}
\Epcp(q)\colon\quad Y^{2}=X(X+1)(X+q^{2}),\qquad q\in\Q^{\ast}\setminus\{0,\pm1\},
\end{equation}
and a rational point $(X,Y)$ is decoded back to a candidate cuboid edge by the
\emph{recovery map}
\begin{equation}\label{eq:phi}
\varphi\colon \Epcp(q)\dashrightarrow\Pone,\qquad
\varphi(X,Y)=\frac{2Yq}{q^{2}-X^{2}}.
\end{equation}
A non-degenerate cuboid requires this value to be finite and nonzero.

The torsion of the family \eqref{eq:Epcp} is $\Z/4\Z\times\Z/2\Z$, uniformly in $q$.
This classification is not new. Yoshida \cite{Yoshida}, studying \emph{face} cuboids,
introduces a family $E_{1,s}\colon y^{2}=x\bigl(x-(2s)^{2}\bigr)\bigl(x+(s^{2}-1)^{2}\bigr)$
(a sub-family of Yagi's curves $y^{2}=x(x-a^{2})(x+b^{2})$ with $a^{2}+b^{2}$ a square)
and proves in his Lemma~2.1 that $E_{1,s}(\Q)_{\mathrm{tor}}\cong\Z/2\Z\times\Z/4\Z$ for
all $s\in\Q\setminus\{0,\pm1\}$, by Mazur's theorem followed by a descent that lands on
$y^{2}=x^{3}-x$. Our family $\Epcp(q)$ obeys the same conclusion by the same route;
the path differs only in bookkeeping, the descent here terminating on the
rank-zero curve $y^{2}=x^{3}+4x$ (Cremona $\mathtt{32a1}$) via the equation
$u^{2}=s^{4}+1$, that is, Fermat's theorem on right triangles. We therefore state the
torsion classification as a recalled proposition and attribute it to Yoshida; the only
presentational addition is an explicit uniform discriminant identity that isolates the
order-eight obstruction directly in the parameter $q$.

It is worth recording, since an earlier draft of the present framework asserted
otherwise, that $\Epcp(q)$ and $E_{1,s}$ are \emph{not} $\Q$-isomorphic on the
Pythagorean locus. Their $j$-invariants coincide only at the non-Pythagorean value
$q=(s^{2}+1)/(2s)$; for genuinely Pythagorean $q=(t^{2}-1)/(2t)$ no rational $s$ matches
the $j$-invariant of $\Epcp(q)$, as a search over $s$ of bounded height confirms. The
two families are related by a $-1$ quadratic twist rather than by isomorphism. The
torsion conclusion survives this distinction because it rests on the rational
two-torsion (three rational roots) and a rational point of order four, features the
twist shares; thus Yoshida's Lemma~2.1 is the correct prior art for the
\emph{classification and its method}, applied here to a twist-inequivalent curve.

The genuinely new content of this note is one statement about the interaction of the
torsion with the recovery map \eqref{eq:phi}, a direction absent from Yoshida's paper,
whose theorems run the opposite way (he uses non-torsion points to \emph{construct}
infinitely many face cuboids). We show that the recovery map is blind to torsion: every
torsion point is degenerate. The three two-torsion points map to the edge value $0$, and
the four order-four points, sitting exactly on the pole divisor $X=\pm q$ of $\varphi$,
map to a pole. Hence no torsion point yields a finite nonzero edge, and any
perfect-cuboid point on the fibre must be of infinite order. The result is elementary
once the torsion structure is in hand, and we make no claim beyond it: it does not
resolve the perfect-cuboid problem, and it is offered as a structural companion to
\cite{Yoshida} and as the foundational lemma for a subsequent rank- or Chabauty-type
analysis of the non-torsion locus.

\section{The curve and its torsion}

We collect the standard facts, all uniform in $q$.

\subsection{Rational two-torsion}
Expanding \eqref{eq:Epcp} gives $Y^{2}=X^{3}+(1+q^{2})X^{2}+q^{2}X$. The right side
factors over $\Q$ as $X(X+1)(X+q^{2})$, so the three nontrivial $2$-torsion points are
\begin{equation}\label{eq:2tors}
T_{1}=(0,0),\qquad T_{2}=(-1,0),\qquad T_{3}=(-q^{2},0),
\end{equation}
giving $(\Z/2\Z)^{2}\subseteq\Epcp(q)(\Q)_{\mathrm{tor}}$ for every
$q\in\Q^{\ast}\setminus\{0,\pm1\}$.

\subsection{Rational four-torsion}
Substituting $X=q$ into \eqref{eq:Epcp} yields $Y^{2}=q(q+1)(q+q^{2})=q^{2}(q+1)^{2}$, so
$Y=\pm q(q+1)$; substituting $X=-q$ yields $Y^{2}=q^{2}(q-1)^{2}$, so $Y=\pm q(q-1)$. The
duplication formula on \eqref{eq:Epcp} gives, with $a_{2}=1+q^{2}$, $a_{4}=q^{2}$ and
$P=(q,q(q+1))$, the slope $\lambda=q+1$ and
\begin{equation}\label{eq:double}
X(2P)=\lambda^{2}-a_{2}-2q=0,\qquad Y(2P)=\lambda\bigl(q-X(2P)\bigr)-q(q+1)=0,
\end{equation}
an identity in $\Q(q)$, so $2P=(0,0)$. The same computation with the sign variants gives
$2\cdot(\pm q,\pm q(q\pm1))=(0,0)$. Hence the four points $(\pm q,\pm q(q\pm1))$ have
order $4$ and
\begin{equation}\label{eq:contain}
\Z/4\Z\times\Z/2\Z\subseteq\Epcp(q)(\Q)_{\mathrm{tor}}
\quad\text{for all }q\in\Q^{\ast}\setminus\{0,\pm1\}.
\end{equation}

\subsection{The classification (after Yoshida)}
The torsion is exactly $\Z/4\Z\times\Z/2\Z$, and this is Yoshida's result transported
to the present curve.

\begin{proposition}[Torsion classification; cf.\ Yoshida {\cite[Lemma 2.1]{Yoshida}}]
\label{prop:tors}
For every Pythagorean $q\in\Q^{\ast}\setminus\{0,\pm1\}$,
\[
\Epcp(q)(\Q)_{\mathrm{tor}}\;=\;\Z/4\Z\times\Z/2\Z,
\]
the eight points being the identity, the three of \eqref{eq:2tors}, and the four
$(\pm q,\pm q(q\pm1))$.
\end{proposition}

\begin{proof}[Proof (recalled)]
Containment is \eqref{eq:contain}. For maximality, Mazur's theorem \cite{Mazur} leaves
only $\Z/2\Z\times\Z/4\Z$, $\Z/2\Z\times\Z/6\Z$, $\Z/2\Z\times\Z/8\Z$ as supergroups of
$\Z/2\Z\times\Z/4\Z$; the middle option fails Lagrange ($8\nmid12$), so it suffices to
exclude a rational point of order $8$. Such a point $P$ has $2P$ of order $4$, hence
$X(2P)\in\{q,-q\}$. Writing $X(2P)=q$ via the duplication formula and clearing
denominators gives the quartic
\[
G_{1}(X,q)=X^{4}-4qX^{3}-(4q^{3}+2q^{2}+4q)X^{2}-4q^{3}X+q^{4},
\]
irreducible over $\Q[X,q]$. Since $G_{1}(0,q)=q^{4}\neq0$, any rational root is
nonzero; setting $X=qZ$ and dividing by $q^{3}$ reduces $G_{1}=0$ to the quadratic in $q$
\[
4Z^{2}q^{2}-\bigl(Z^{4}-4Z^{3}-2Z^{2}-4Z+1\bigr)q+4Z^{2}=0,
\]
whose discriminant is
\begin{equation}\label{eq:disc}
\Delta(Z)=\bigl(Z^{4}-4Z^{3}-2Z^{2}-4Z+1\bigr)^{2}-64Z^{4}
=(Z-1)^{4}(Z+1)^{2}(Z^{2}-6Z+1).
\end{equation}
The prefactor $(Z-1)^{4}(Z+1)^{2}=\bigl((Z-1)^{2}(Z+1)\bigr)^{2}$ is a square, so
$\Delta(Z)$ is a square exactly when $Z^{2}-6Z+1=w^{2}$. This conic has the rational
point $(0,1)$; parametrising by $w=tZ+1$ gives $Z=Z(t)=-2(t+3)/\bigl((t-1)(t+1)\bigr)$,
and back-substitution simplifies the recovered $q$ to
\[
q=q_{+}(t)=\Bigl(\tfrac{4(t+1)}{(t-1)(t+3)}\Bigr)^{2}=s^{2},
\qquad q_{-}(t)=1/s^{2}.
\]
Thus any order-eight parameter forces $q=s^{2}$ or $1/s^{2}$. The Pythagorean hypothesis
$1+q^{2}\in\Q^{2}$ then demands $s^{4}+1\in\Q^{2}$, i.e.\ a rational point on
$u^{2}=s^{4}+1$. This curve is birational to $y^{2}=x^{3}+4x$ (conductor $32$, Cremona
$\mathtt{32a1}$), of analytic rank $0$ and torsion $\Z/4\Z$; by Kolyvagin its
Mordell--Weil rank is $0$, so the only rational solution is $s=0$, giving
$q\in\{0,\infty\}$, both excluded. The case $X(2P)=-q$ is identical with $q\mapsto-q$.
The reduction to $s^{4}+1=u^{2}$ and its resolution by Fermat's theorem on right
triangles is exactly the descent of Yoshida \cite[Lemma 2.1]{Yoshida}, there phrased on
$y^{2}=x^{3}-x$; the discriminant identity \eqref{eq:disc} is the only presentational
change.
\end{proof}

\begin{remark}
Equation \eqref{eq:disc} is the single point at which the present write-up departs from
\cite{Yoshida} at the level of formulae: it exhibits the order-eight obstruction as a
uniform discriminant in the parameter $q$, with the square prefactor stripped off,
rather than passing through an auxiliary curve. It is a convenience, not a new theorem.
\end{remark}

\section{Degeneracy of torsion under the recovery map}

The recovery map \eqref{eq:phi} is regular on the affine locus $X\neq\pm q$ and has pole
divisor $\{X=q\}\cup\{X=-q\}$. We record its action on torsion. This is the new content
of the note.

\begin{theorem}[Torsion degeneracy]\label{thm:main}
Let $q\in\Q^{\ast}\setminus\{0,\pm1\}$ be Pythagorean. Then every one of the seven
non-identity rational torsion points of $\Epcp(q)$ is sent by $\varphi$ into
$\{0,\infty\}$: the three $2$-torsion points $\eqref{eq:2tors}$ map to $\varphi=0$, and
the four order-four points $(\pm q,\pm q(q\pm1))$ lie on the pole divisor of $\varphi$,
hence map to $\infty$. (The identity $O$ is the point at infinity of $\Epcp(q)$, outside
the affine domain of $\varphi$.) Consequently no rational torsion point of $\Epcp(q)$ has
finite nonzero image under $\varphi$.
\end{theorem}

\begin{proof}
By Proposition~\ref{prop:tors} the torsion consists of the identity, the three points of
\eqref{eq:2tors}, and the four order-four points.

For a $2$-torsion point $T_{i}=(X_{i},0)$ the numerator $2Yq$ of \eqref{eq:phi} vanishes,
while the denominator $q^{2}-X_{i}^{2}$ does not: for $T_{1}$ it is $q^{2}\neq0$, for
$T_{2}$ it is $q^{2}-1\neq0$ (as $q\neq\pm1$), and for $T_{3}$ it is
$q^{2}-q^{4}=q^{2}(1-q^{2})\neq0$. Hence $\varphi(T_{i})=0$ for $i=1,2,3$.

For an order-four point $(\pm q,\pm q(q\pm1))$ the $X$-coordinate satisfies $X^{2}=q^{2}$,
so the denominator $q^{2}-X^{2}$ is $0$, while the numerator $2Yq=\pm2q^{2}(q\pm1)$ is
nonzero because $q\notin\{0,\pm1\}$. Each of the four points therefore lies on the pole
divisor, and $\varphi$ takes the value $\infty$ there.

The identity $O$ is the point at infinity of $\Epcp(q)$, outside the affine domain of
$\varphi$, and so is not assigned a finite nonzero edge; it is excluded from the seven
counted above. The two displayed evaluations are identities in $\Q(q)$, independent of
the particular Pythagorean value of $q$.
\end{proof}

\begin{corollary}\label{cor:nontorsion}
On any Pythagorean fibre, a rational point of $\Epcp(q)$ decoding to a finite nonzero
cuboid edge is necessarily of infinite order. In particular the torsion subgroup may be
discarded at the outset of any search for perfect-cuboid solutions on the fibre.
\end{corollary}

\begin{proof}
A finite nonzero value of $\varphi$ excludes, by Theorem~\ref{thm:main}, every torsion
point. Since $\Epcp(q)(\Q)=\Epcp(q)(\Q)_{\mathrm{tor}}\oplus\Z^{r}$ with the torsion
exhausted by Proposition~\ref{prop:tors}, such a point has infinite order.
\end{proof}

\begin{remark}
Yoshida \cite{Yoshida} has no analogue of Theorem~\ref{thm:main}: his map runs from
non-torsion points to face cuboids, producing solutions, whereas $\varphi$ here detects
the \emph{absence} of solutions among torsion. The two viewpoints are compatible -- a
perfect cuboid is a face cuboid with one further squareness condition -- and the present
statement is the structural reason the torsion of \eqref{eq:Epcp} contributes nothing to
the perfect-cuboid count.
\end{remark}

\section{Computational verification}

All claims were checked symbolically in \texttt{sympy} and numerically in PARI/GP; the
scripts and their output accompany this note.

\begin{enumerate}
\item \emph{Discriminant identity \eqref{eq:disc}.} The reduction $X(2P)=q\Rightarrow
G_{1}=0$, the substitution $X=qZ$, and the factorisation
$\Delta(Z)=(Z-1)^{4}(Z+1)^{2}(Z^{2}-6Z+1)$ were verified as exact identities, together
with the conic parametrisation giving $q=s^{2}$
(\texttt{03\_discriminant\_identity.py}).
\item \emph{Degeneracy (Theorem~\ref{thm:main}).} The action of $\varphi$ on the eight
torsion points was confirmed as an identity in $\Q(q)$
(\texttt{05\_recovery\_degeneracy.py}) and numerically over $62$ Pythagorean values of
$q$, all $62\times8=496$ evaluations landing in $\{0,\infty\}$ with torsion $[4,2]$ in
every case (\texttt{02\_torsion\_sweep.gp}).
\item \emph{Fermat obstruction.} The curve $y^{2}=x^{3}+4x$ has analytic rank $0$ and
torsion $\Z/4\Z$, so $u^{2}=s^{4}+1$ has only $s=0$ (\texttt{04\_fermat\_curve.gp}).
\item \emph{Twist, not isomorphism.} For Pythagorean $q\in\{4/3,3/4,5/12,12/5,15/8,8/15\}$
no rational $s$ of height $\le80$ has $j(E_{1,s})=j(\Epcp(q))$, whereas the
$j$-invariants agree along the non-Pythagorean locus $q=(s^{2}+1)/(2s)$; the curves are
$-1$-twist-related (\texttt{01\_isomorphism.gp}).
\end{enumerate}

\section{Scope}

This is a short structural note. Its single new theorem, the torsion degeneracy
\ref{thm:main}, is a companion observation to Yoshida \cite{Yoshida}: it identifies the
recovery map under which the entire torsion subgroup of the perfect-cuboid family is
degenerate, and thereby justifies restricting any subsequent rank-, descent-, or
Chabauty-type search to the non-torsion locus (Corollary~\ref{cor:nontorsion}). The
torsion classification itself is Yoshida's and is presented as such. We make no claim
about the existence or non-existence of perfect cuboids; the content here is the
elementary but, to our knowledge, previously unrecorded fact that torsion alone never
produces a finite nonzero edge.

\section*{Acknowledgements}
The author thanks the maintainers of PARI/GP and SymPy.

\begin{thebibliography}{9}

\bibitem{Guy} R.~K.~Guy, \emph{Unsolved Problems in Number Theory}, 3rd ed.,
Problem Books in Mathematics, Springer, New York, 2004.

\bibitem{Mazur} B.~Mazur, \emph{Rational isogenies of prime degree}, Invent. Math.
\textbf{44} (1978), 129--162.

\bibitem{Yoshida} T.~Yoshida, \emph{The relationship between face cuboids and elliptic
curves}, preprint, arXiv:2407.09825 [math.NT], 2024.

\end{thebibliography}

\end{document}
